\documentclass[10pt,a4paper]{amsart}
\usepackage[margin=19mm]{geometry}
\usepackage{amsmath,amssymb,mathtools}
\usepackage[scr=rsfs,cal=euler]{mathalfa}
\usepackage{xfrac}
\usepackage[unicode,hidelinks,bookmarks=false]{hyperref}
\newcommand{\ie}{i.e.\ }
\newcommand{\cf}{cf.\ }
\newcommand{\resp}{respectively\ }
\newcommand{\Subsection}{\subsection}
\providecommand{\nobreakdash}{\nobreak\hbox{-}}
\setlength{\emergencystretch}{2em}
\multlinegap0pt
\newtheorem{asmp}{Asmp}
\newtheorem{lemma}{Lemma}
\title{Capital Asset Pricing Model with Size Factor and Normalizing by Volatility Index}
\author{Abraham Atsiwo, Andrey Sarantsev}
\date{Source: arXiv:2411.19444v8 (CC BY 4.0)}
\begin{document}
\maketitle

\noindent\emph{Source and licence.} Capital Asset Pricing Model with Size Factor and Normalizing by Volatility Index. Version arXiv:2411.19444v8 (CC BY 4.0), distributed by arXiv under the Creative Commons Attribution 4.0 International licence, http://creativecommons.org/licenses/by/4.0/, as recorded in the arXiv licence metadata for this exact version and shown on the abstract page https://arxiv.org/abs/2411.19444v8. Full licence notice: \texttt{licenses/arxiv-2411.19444-CC-BY-4.0.txt}.\par\smallskip

\noindent\textit{Editorial scope.} Counted unit: the article's lemma lemma (source0.tex lines 312--313). The article's complete original proof of it is delivered with it (source0.tex lines 315--328, one \texttt{proof} environment of 14 lines). No mathematics is written, repaired or re-derived: the statement and the proof are the article's own lines, in the article's own notation, and no retained line is substituted or re-flowed.  Retained uncounted as the source-local context the proof needs: lines 147--150; lines 271--274; lines 277--280; lines 282--285; lines 287--290; lines 292--295; lines 297--300.  Nothing was omitted from any retained span, so the package carries no omission record.  No retained line contains a citation, so the wrapper carries no bibliography and no bibliography entry.  No retained line contains a figure environment, an image-inclusion command or drawing code, so no image is delivered and the package declares no asset dependency.  Editorial formatting only, in addition: an AMS \texttt{amsart} preamble that restates the article's own declarations and macros used by the retained lines, namely the environments \texttt{asmp}, \texttt{lemma} on separate counters, together with the standard packages those lines need.  The retained environments print in this document as Asmp 1, Lemma 1 in order of appearance, whereas the article prints the same items with the numbering of its own full document.  The complete native source of the article as distributed in the arXiv e-print for this exact version is bundled unchanged as \texttt{sources/169320/source0.tex}.
\begin{equation}
\label{eq:ln-VIX}
\ln V(t) = \alpha + \beta \ln V(t-1) + W(t).
\end{equation}

\begin{asmp}
Vectors $\mathbf{Y}(t)$ are IID with mean zero and finite second moment, with Lebesgue joint density on $\mathbb R^5$ which is everywhere strictly positive. 
\label{asmp:innovations}
\end{asmp}

\begin{equation}
\label{eq:n-return}
\frac{R_0(t)}{V(t)} = g + U(t),
\end{equation}

\begin{equation}
\label{eq:n-premia}
\frac{P_0(t)}{V(t)} = G + Z(t),
\end{equation}

\begin{equation}
\label{eq:premia-main}
\frac{P(t)}{V(t)} = M + AC(t) + (1 + BC(t))\frac{P_0(t)}{V(t)} + \delta(t).
\end{equation}

\begin{equation}
\label{eq:return-main}
\frac{R(t)}{V(t)} = m + aC(t) + (1 + bC(t))\frac{R_0(t)}{V(t)} + \varepsilon(t).
\end{equation}

\begin{equation}
\label{eq:rel-size}
C(t) = \ln\frac{S(t)}{S_0(t)}.
\end{equation} 

\begin{lemma} Under Assumption~\ref{asmp:innovations}, the process $C$ is Markov. Also, the truncated model $(V, R_0, R)$ is Markov. Finally, the completed model $(V, R_0, R, P_0, P)$ is Markov. 
\end{lemma}

\begin{proof} It is clear from~\eqref{eq:ln-VIX} that the process $\ln V$ (and therefore $V$) is Markov. Together with~\eqref{eq:n-return}, this shows that $(V, R_0)$ is Markov. From definition~\eqref{eq:rel-size}, we write 
\begin{align}
\label{eq:dc}
\begin{split}
C(t+1) - C(t) &= \ln\frac{S(t+1)}{S_0(t+1)} - \ln\frac{S(t)}{S_0(t)} \\ & = \ln\frac{S(t+1)}{S(t)} - \ln\frac{S_0(t+1)}{S_0(t)} = R(t) - R_0(t).
\end{split}
\end{align}
Using~\eqref{eq:dc}, we rearrange~\eqref{eq:return-main} as follows:
\begin{equation}
\label{eq:C-main}
C(t+1) = (1 + aV(t) + bR_0(t))C(t) + V(t)(m + \varepsilon(t)).
\end{equation}
But from Assumption~\ref{asmp:innovations}, this process $C$ from~\eqref{eq:C-main} is Markov, too. From~\eqref{eq:dc} and~\eqref{eq:n-return}, this equation~\eqref{eq:C-main} shows that $(V, R_0, C)$, or, equivalently, $(V, R_0, R)$, is Markov. Finally, from~\eqref{eq:premia-main} and~\eqref{eq:n-premia}, we get that the completed model is Markov. 
\end{proof}

\end{document}
